\documentclass[presentation,aspectratio=169]{ctexbeamer}

\setbeamertemplate{sidebar left}[frame number]
\hypersetup{pdfpagemode=FullScreen}
\usetheme{SCUT}

%----------------------------------------------------------------------------------------
%	TITLE PAGE
%----------------------------------------------------------------------------------------

\title{代码覆盖率及路径分析检测}
\subtitle{软件测试课程大作业}
\author{徐宇鸿}
\institute
{
计算机科学与工程学院 \\ % Your institution for the title page
\medskip
}
\date{\today}

\begin{document}

\begin{frame}
\titlepage
\end{frame}

\begin{frame}
\frametitle{目录}
\tableofcontents
\end{frame}

%========================================================================================
%	第一部分：什么是路径分析
%========================================================================================

\section{什么是路径分析与覆盖率}

\begin{frame}[fragile]
\frametitle{程序执行路径与三种基本结构}
{\small 程序运行时并非所有代码都会被执行——\textbf{条件分支}和\textbf{循环}决定了程序走哪条路。%
\textbf{路径分析}就是追踪程序\emph{实际走了哪些路}，找出\emph{哪些路从未走过}。}

\vspace{0.15cm}
\begin{columns}[T]
\begin{column}{0.30\textwidth}
\begin{block}{\small 顺序结构}
{\scriptsize
代码从上到下依次执行，\textbf{只有一条路径}，不存在分支。\\[2pt]
\texttt{x = 1}\\
\texttt{y = x + 2}\\
\texttt{print(y)}}
\end{block}
\end{column}
\begin{column}{0.32\textwidth}
\begin{block}{\small 选择结构（分支）}
{\scriptsize
根据条件走\textbf{不同分支}，产生\textbf{多条路径}。\\[2pt]
\texttt{if x > 0:}\\
\texttt{~~return "正"}\\
\texttt{elif x == 0:}\\
\texttt{~~return "零"}\\
\texttt{else:}\\
\texttt{~~return "负"}}
\end{block}
\end{column}
\begin{column}{0.34\textwidth}
\begin{block}{\small 循环结构}
{\scriptsize
重复执行某段代码，循环次数影响路径。\\[2pt]
\texttt{for i in range(n):}\\
\texttt{~~total += i}\\[3pt]
\texttt{while x > 0:}\\
\texttt{~~x = x // 2}}
\end{block}
\end{column}
\end{columns}

\vspace{0.15cm}
{\small\textbf{分支覆盖率}衡量的是：程序中所有分支点里，有多少被测试执行过。}
\[
\text{分支覆盖率} = \frac{\text{已执行的分支数}}{\text{总分支数}} \times 100\%
\]

{\small 例：程序有 5 个分支（3 个 \texttt{if} + 1 个 \texttt{for} + 1 个 \texttt{try}），运行后走了 3 个 $\Rightarrow$ 覆盖率 $60\%$。覆盖率越高 $\Rightarrow$ 测试越充分，隐藏 bug 越少。}
\end{frame}

%------------------------------------------------

\begin{frame}
\frametitle{路径分析与软件测试的关系}
{\small 路径分析是白盒测试的核心手段——通过追踪程序内部执行路径，量化测试的\textbf{充分性}。}

\vspace{0.15cm}
\begin{columns}[T]
\begin{column}{0.55\textwidth}
{\small\textbf{测试流程：}}
\begin{enumerate}\setlength{\itemsep}{1pt}
    \item \textbf{静态分析}：解析源码，找出所有分支点（\texttt{if/for/while/try}）
    \item \textbf{设计测试用例}：构造不同输入，尽可能覆盖每条分支
    \item \textbf{动态追踪}：运行程序，记录实际走了哪些行/分支
    \item \textbf{计算覆盖率}：对比静态分支与动态执行结果
    \item \textbf{补充测试}：对未覆盖分支补充测试用例，提升覆盖率
\end{enumerate}
\end{column}
\begin{column}{0.42\textwidth}
\begin{alertblock}{\small 为什么要做路径分析？}
{\scriptsize
如果不追踪执行路径，你\textbf{无法知道}某段 \texttt{else} 分支是否被测试过。\\[3pt]
\textbf{未覆盖 $=$ 未测试 $=$ 可能藏有 bug}\\[3pt]
典型场景：
\begin{itemize}\setlength{\itemsep}{0pt}
    \item 异常处理（\texttt{except}）从未触发
    \item 边界条件（\texttt{elif}）从未命中
    \item 循环体从未进入或从未跳过
\end{itemize}}
\end{alertblock}
\end{column}
\end{columns}

\vspace{0.15cm}
{\small\textbf{本作业目标}：编写 PathTracer 工具，自动完成上述步骤 1–4，生成描述程序执行分支的覆盖率报告。}
\end{frame}

%========================================================================================
%	第二部分：核心概念
%========================================================================================

\section{核心概念与工作原理}

\begin{frame}
\frametitle{整体流程概览}
路径分析工具的工作分为三个阶段：

\vspace{0.5cm}
\begin{enumerate}
    \item \textbf{静态分析}——不运行程序，用 Python 的 \texttt{ast} 模块解析源码，找出所有分支点（if / for / while / try）
    \vspace{0.3cm}
    \item \textbf{动态追踪}——运行程序，用 \texttt{sys.settrace()} 在虚拟机级别钩住每一步执行，记录实际走了哪些行
    \vspace{0.3cm}
    \item \textbf{覆盖率计算}——对比两份结果：分支点所在行是否出现在已执行行号集合中，得到覆盖率百分比
\end{enumerate}
\end{frame}

%------------------------------------------------

\begin{frame}[fragile]
\frametitle{静态分析（1/2）：AST——抽象语法树是什么？}
{\small Python 执行代码前会先把源码经过\textbf{词法分析}$\rightarrow$\textbf{语法分析}$\rightarrow$变成 AST（一棵树），再由虚拟机执行。%
\texttt{ast} 模块让我们能直接访问这棵树。}

\vspace{0.2cm}
\begin{columns}[T]
\begin{column}{0.45\textwidth}
\begin{block}{\small 一段简单的源码}
{\small
\begin{verbatim}
def f(x):
    if x > 0:
        return 1
    else:
        return -1
\end{verbatim}}
\end{block}
\end{column}
\begin{column}{0.52\textwidth}
\begin{block}{\small 对应的 AST 树形结构}
{\small
\texttt{Module} \\
\texttt{~~$\hookrightarrow$ FunctionDef(name='f')} \\
\texttt{~~~~~~$\hookrightarrow$ If(test=Compare(x>0))} \\
\texttt{~~~~~~~~~~$\hookrightarrow$ Return(Constant(1))} \\
\texttt{~~~~~~~~~~$\hookrightarrow$ orelse:} \\
\texttt{~~~~~~~~~~~~~~$\hookrightarrow$ Return(Constant(-1))}}

\vspace{0.15cm}
{\small 每个 \texttt{if}/\texttt{for}/\texttt{while}/\texttt{try} 在树中都有\textbf{对应的节点类型}：\texttt{ast.If}、\texttt{ast.For}……}
\end{block}
\end{column}
\end{columns}

\vspace{0.2cm}
{\small\textbf{关键点：}我们不需要自己写解析器——Python 标准库 \texttt{ast.parse()} 一行调用就能把任意源码变成这棵树。然后用 \texttt{ast.walk(tree)} 深度优先遍历所有节点，逐个检查类型即可。}
\end{frame}

%------------------------------------------------

\begin{frame}[fragile]
\frametitle{静态分析（2/2）：项目中的实现——\texttt{cfg\_builder.py}}
{\small 前一页讲了 AST 的原理。现在来看我们项目中\textbf{实际写的代码}：\\
\texttt{cfg\_builder.py} 用 \texttt{ast} 模块遍历语法树，识别每种控制流节点，提取分支信息。}

\vspace{0.05cm}
{\scriptsize
\begin{verbatim}
class CFGBuilder:
    def build_from_file(self, file_path):
        source = open(file_path).read()
        tree = ast.parse(source)               # 源码 -> AST树
        for node in ast.walk(tree):            # 遍历所有节点
            branch = self._extract_branch(node)
            if branch:
                self.branches.append(branch)

    def _extract_branch(self, node):
        if isinstance(node, ast.If):
            return CodeBranch(line_no=node.lineno,
                branch_type=BranchType.IF, end_line=node.end_lineno)
        elif isinstance(node, ast.For):
            return CodeBranch(..., BranchType.FOR)
        elif isinstance(node, ast.While):
            return CodeBranch(..., BranchType.WHILE)
        elif isinstance(node, ast.ExceptHandler):
            return CodeBranch(..., BranchType.EXCEPT)
        return None   # 其他节点不产生分支
\end{verbatim}}
\end{frame}

%------------------------------------------------

\begin{frame}[fragile]
\frametitle{前置知识（1/2）：源码 $\rightarrow$ 字节码}
{\footnotesize 要理解动态追踪的原理，需要先了解 Python 代码的执行过程。}

\vspace{0.1cm}
\begin{block}{\footnotesize Python \textbf{不会}直接执行你写的代码}
{\footnotesize 它先把源码\textbf{编译}成``字节码''（Bytecode）——一种简化的中间指令。源码的一行可能对应多条字节码指令。\\
可以用 Python 内置的 \texttt{dis} 模块查看：}
\end{block}

\vspace{0.1cm}
\begin{columns}[T]
\begin{column}{0.35\textwidth}
\begin{block}{\footnotesize 你写的源码}
{\footnotesize
\begin{verbatim}
x = 1 + 2
\end{verbatim}}
\end{block}
\end{column}
\begin{column}{0.60\textwidth}
\begin{block}{\footnotesize 翻译后的字节码}
{\footnotesize
\begin{verbatim}
LOAD_CONST    1    # 取出数字 1
LOAD_CONST    2    # 取出数字 2
BINARY_ADD         # 做加法
STORE_NAME    x    # 存到变量 x
\end{verbatim}}
\end{block}
\end{column}
\end{columns}

\vspace{0.1cm}
{\footnotesize \textbf{类比：}源码像``中文''，字节码像``英文''，虚拟机能读懂英文并逐条执行。我们不需要手写字节码——Python 自动完成翻译。}
\end{frame}

%------------------------------------------------

\begin{frame}
\frametitle{前置知识（2/2）：虚拟机与栈帧}
{\footnotesize CPython 虚拟机（Virtual Machine）是一个用 C 语言写的\textbf{大循环}，它逐条读取字节码并执行。}

\vspace{0.1cm}
\begin{block}{\footnotesize 栈帧（Frame）——虚拟机的``笔记本''}
{\footnotesize 每条字节码执行时，虚拟机维护一个\textbf{栈帧}，里面记录着当前执行状态：\\[3pt]
\quad $\bullet$ \texttt{f\_lineno} —— 当前执行到源码\textbf{第几行}\\
\quad $\bullet$ \texttt{f\_locals} —— 当前所有\textbf{局部变量}的值\\
\quad $\bullet$ \texttt{f\_code.co\_name} —— 当前在\textbf{哪个函数}里\\
\quad $\bullet$ \texttt{f\_code.co\_varnames} —— 函数的\textbf{参数名}列表}
\end{block}

\vspace{0.1cm}
\begin{block}{\footnotesize 调用栈（Call Stack）——函数嵌套的``俄罗斯套娃''}
{\footnotesize 当函数 A 调用函数 B 时，虚拟机会创建一个新的栈帧压入\textbf{调用栈}。函数返回时弹出。\\[3pt]
调用栈让虚拟机知道：``执行完后该回到哪里继续''。\\[3pt]
\textbf{关键：}\texttt{sys.settrace()} 就是让虚拟机在每条字节码执行前后``暂停一下''，把当前栈帧交给你检查。这样你就能知道程序执行到了哪一行、变量是什么值。}
\end{block}
\end{frame}

%------------------------------------------------

\begin{frame}
\frametitle{动态追踪（1/2）：\texttt{sys.settrace()} 的工作机制}
{\footnotesize \texttt{sys.settrace()} 是 CPython 虚拟机提供的\textbf{调试钩子}——不需要修改被追踪程序的源码，从``外部''监控执行。}

\vspace{0.1cm}
\begin{block}{\footnotesize 底层原理}
{\footnotesize CPython 的主执行循环（\texttt{Python/ceval.c}）在每条\textbf{字节码}执行前后，都会检查是否设置了 trace function（即你通过 \texttt{sys.settrace(回调函数)} 注册的那个函数）。如果有，就暂停执行，把\textbf{栈帧}传给你的回调函数。这是\textbf{解释器级别}的钩子——不需要修改被追踪的代码。}
\end{block}

\vspace{0.1cm}
\begin{block}{\footnotesize 回调函数接收的三个参数}
{\footnotesize
\begin{tabular}{lp{7.5cm}}
\toprule
\textbf{参数} & \textbf{含义} \\
\midrule
\texttt{frame} & 当前\textbf{栈帧对象}：\texttt{f\_lineno}（行号）、\texttt{f\_code.co\_name}（函数名）、\texttt{f\_locals}（局部变量） \\
\texttt{event} & 事件类型：\texttt{'call'}（函数调用）、\texttt{'line'}（执行一行）、\texttt{'return'}（返回）、\texttt{'exception'}（异常） \\
\texttt{arg} & 附加数据：\texttt{'return'} $\rightarrow$ 返回值；\texttt{'exception'} $\rightarrow$ \texttt{(exc\_type, exc\_value, traceback)} \\
\bottomrule
\end{tabular}}
\end{block}
\end{frame}

%------------------------------------------------

\begin{frame}[fragile]
\frametitle{动态追踪（2/2）：项目中的实现——\texttt{tracer.py}}
{\footnotesize 前一页讲了 \texttt{sys.settrace()} 的原理。现在来看我们项目中\textbf{实际写的代码}：\\
\texttt{tracer.py} 里的 \texttt{\_trace\_callback} 方法就是注册给 \texttt{sys.settrace()} 的那个回调函数。\\
它根据不同事件分别处理：记录行号、函数调用链和返回值。}

\vspace{0.05cm}
{\tiny
\begin{verbatim}
def _trace_callback(self, frame, event, arg):
    if not frame.f_code.co_filename.endswith(self._target_file):
        return self._trace_callback         # 只追踪目标文件

    if event == 'line':
        self.executed_lines.add(frame.f_lineno)             # 记录行号

    elif event == 'call':
        args = {k: repr(frame.f_locals[k])                  # 提取参数
                for k in frame.f_code.co_varnames}
        call = FunctionCall(function_name=frame.f_code.co_name,
                            call_depth=len(self._call_stack), arguments=args)
        if self._call_stack:
            self._call_stack[-1].children.append(call)      # 子调用
        else:
            self.function_calls.append(call)                 # 顶层调用
        self._call_stack.append(call)                        # 入栈

    elif event == 'return':
        call = self._call_stack.pop()                        # 出栈
        call.return_value = repr(arg)                        # 记录返回值
\end{verbatim}}

\vspace{0.05cm}
{\footnotesize\textbf{要点：}\texttt{\_call\_stack} 列表模拟函数调用栈 $\rightarrow$ 还原完整调用树（谁调用谁、参数、返回值）。}
\end{frame}

%------------------------------------------------

\begin{frame}
\frametitle{覆盖率计算与报告生成}
{\small \texttt{reporter.py} 将静态分析的分支列表与动态追踪的行号集合\textbf{求交集}，得到覆盖情况。}

\vspace{0.1cm}
\begin{columns}[T]
\begin{column}{0.55\textwidth}
\begin{block}{\small 判定逻辑（逐分支检查）}
{\scriptsize
\begin{enumerate}\setlength{\itemsep}{1pt}
    \item \texttt{cfg\_builder} 解析源码 $\rightarrow$ 得到所有 \texttt{CodeBranch} 列表，每个记录 \texttt{line\_no} 和 \texttt{branch\_type}
    \item \texttt{tracer} 运行程序 $\rightarrow$ 得到 \texttt{executed\_lines: Set[int]}（已执行行号集合）
    \item 对每个分支：\\
    \quad \texttt{if branch.line\_no in executed\_lines:}\\
    \quad\quad \texttt{branch.is\_executed = True}
    \item 统计覆盖率：
    \[
        \text{coverage} = \frac{\sum \mathbf{1}[\text{branch.is\_executed}]}{|\text{branches}|} \times 100\%
    \]
    \item 按类型分组统计：if / for / while / try / except 各多少
\end{enumerate}}
\end{block}
\end{column}
\begin{column}{0.42\textwidth}
\begin{exampleblock}{\small 具体示例}
{\scriptsize
\textbf{静态分析结果（10 个分支）：}\\
Line 2: \texttt{if} \checkmark~Line 4: \texttt{if} \checkmark~Line 9: \texttt{for} \checkmark\\
Line 14: \texttt{if} \checkmark~Line 18: \texttt{while} \checkmark~Line 22: \texttt{try} \checkmark\\
Line 25: \texttt{except} \textcolor{red}{\texttimes}~Line 28: \texttt{for} \checkmark\\
Line 32: \texttt{if} \textcolor{red}{\texttimes}~Line 35: \texttt{while} \textcolor{red}{\texttimes}\\[3pt]
\textbf{已执行 7/10，覆盖率 70\%}\\[3pt]
未覆盖：\texttt{except}（没触发异常）、\texttt{if}（条件未命中）、\texttt{while}（循环未进入）\\[3pt]
\textbf{启示：}需要构造能触发异常、命中特定条件、进入循环的测试用例。}
\end{exampleblock}
\end{column}
\end{columns}
\end{frame}

%========================================================================================
%	第三部分：PathTracer 工具实现
%========================================================================================

\section{PathTracer 工具实现}

\begin{frame}
\frametitle{项目模块结构}
\begin{columns}[T]
\begin{column}{0.45\textwidth}
\begin{block}{\footnotesize 文件列表}
{\footnotesize
\begin{tabular}{ll}
\toprule
\textbf{文件} & \textbf{职责} \\
\midrule
\texttt{models.py} & 数据模型定义 \\
\texttt{cfg\_builder.py} & 静态分析（CFG） \\
\texttt{tracer.py} & 动态追踪 \\
\texttt{reporter.py} & 覆盖率计算与报告 \\
\texttt{cli.py} & 命令行入口 \\
\texttt{demo.py} & 演示脚本 \\
\bottomrule
\end{tabular}}
\end{block}
\end{column}

\begin{column}{0.50\textwidth}
\begin{block}{\footnotesize 模块协作流程}
{\footnotesize
\centering
\begin{enumerate}
    \item \texttt{cfg\_builder} 解析源代码
    \item[] $\downarrow$ 得到所有分支点
    \item \texttt{tracer} 追踪运行过程
    \item[] $\downarrow$ 得到已执行行号
    \item \texttt{reporter} 对比 $\rightarrow$ 覆盖率
    \item[] $\downarrow$
    \item 输出报告（文本/HTML/JSON）
\end{enumerate}}
\end{block}
\end{column}
\end{columns}
\end{frame}

%------------------------------------------------

\begin{frame}[fragile]
\frametitle{数据模型（1/2）：分支类型与分支点}
{\footnotesize \texttt{models.py} 定义了整个工具使用的数据结构。涉及两个 Python 基础概念：}

\vspace{0.1cm}
\begin{columns}[T]
\begin{column}{0.55\textwidth}
{\scriptsize
\begin{verbatim}
class BranchType(Enum):
    """所有可识别的分支类型"""
    SEQUENTIAL = "sequential"
    IF     = "if"
    ELSE   = "else"
    FOR    = "for"
    WHILE  = "while"
    TRY    = "try"
    EXCEPT = "except"

@dataclass
class CodeBranch:
    """表示代码中的一个分支点"""
    file_path: str
    line_no: int
    branch_type: BranchType
    end_line: int = 0
    is_executed: bool = False
\end{verbatim}}
\end{column}
\begin{column}{0.42\textwidth}
\begin{block}{\scriptsize \texttt{Enum} 枚举是什么？}
{\scriptsize 枚举就是``一组有名字的常量''。\\[2pt]
用数字记分支类型容易混淆（1=if? 2=for?），用 \texttt{BranchType.IF} 这样的名字就清晰多了。\\[2pt]
Python 标准库 \texttt{enum.Enum} 提供这个功能。}
\end{block}
\vspace{0.05cm}
\begin{block}{\scriptsize \texttt{@dataclass} 装饰器是什么？}
{\scriptsize 普通 class 要自己写 \texttt{\_\_init\_\_} 方法来初始化属性。\texttt{@dataclass} 让 Python \textbf{自动生成}这些代码，你只需声明字段即可。\\[2pt]
相当于自动帮你生成了：\\
\texttt{def \_\_init\_\_(self, file\_path, line\_no, ...):}\\
\texttt{~~self.file\_path = file\_path}\\
\texttt{~~self.line\_no = line\_no}\\
\texttt{~~...}}
\end{block}
\end{column}
\end{columns}
\end{frame}
\end{frame}

%------------------------------------------------

\begin{frame}[fragile]
\frametitle{数据模型（2/2）：函数调用与执行路径}
{\footnotesize 还需要记录\textbf{函数调用链}和\textbf{整体执行路径}：}

\vspace{0.05cm}
{\tiny
\begin{verbatim}
@dataclass
class FunctionCall:
    """记录一次函数调用"""
    function_name: str                       # 函数名
    file_path: str                           # 所在文件
    line_no: int                             # 调用位置行号
    call_depth: int                          # 调用深度（0=顶层调用）
    arguments: Dict[str, str] = field(...)   # 参数名 -> 参数值（字符串）
    return_value: Optional[str] = None       # 返回值（字符串）
    children: List[FunctionCall] = field(...) # 子调用列表（嵌套结构）

@dataclass
class ExecutionPath:
    """记录一个文件的完整执行路径"""
    file_path: str
    executed_lines: Set[int] = field(...)               # 执行过的所有行号
    executed_branches: Dict[int, CodeBranch] = field(...)  # 被执行的分支
    function_calls: List[FunctionCall] = field(...)     # 函数调用链

@dataclass
class CoverageReport:
    """最终的覆盖率报告"""
    total_branches: int = 0
    executed_branches: int = 0
    coverage_percentage: float = 0.0
    branches_by_type: Dict[BranchType, int] = field(...)  # 按类型统计
    file_reports: Dict[str, ExecutionPath] = field(...)   # 按文件统计
\end{verbatim}}
\end{frame}

%========================================================================================
%	第四部分：使用演示
%========================================================================================

\section{使用演示}

\begin{frame}[fragile]
\frametitle{被追踪的示例代码（1/2）：条件分支}
{\footnotesize \texttt{demo.py} 包含多种程序结构，演示 PathTracer 的追踪能力：}

\vspace{0.1cm}
\begin{columns}[T]
\begin{column}{0.48\textwidth}
\begin{block}{\scriptsize if/elif/else 选择结构}
{\scriptsize
\begin{verbatim}
def demo_conditionals(value):
    if value < 0:
        return "negative"
    elif value == 0:
        return "zero"
    elif value <= 10:
        return "small"
    elif value <= 50:
        return "medium"
    else:
        return "large"
\end{verbatim}}
\end{block}
\end{column}
\begin{column}{0.48\textwidth}
\begin{block}{\scriptsize 说明}
{\scriptsize
\textbf{5 条可能的执行路径}，输入不同走的路不同：\\[3pt]
$\bullet$ \texttt{demo\_conditionals(-5)}\\
~~走第1个 if $\rightarrow$ "negative"\\[2pt]
$\bullet$ \texttt{demo\_conditionals(25)}\\
~~走第4个 elif $\rightarrow$ "medium"\\[2pt]
$\bullet$ \texttt{demo\_conditionals(100)}\\
~~走 else $\rightarrow$ "large"\\[3pt]
单一输入\textbf{只能覆盖一条路径}。}
\end{block}
\end{column}
\end{columns}
\end{frame}

%------------------------------------------------

\begin{frame}[fragile]
\frametitle{被追踪的示例代码（2/2）：循环、异常与嵌套}
{\footnotesize 除了 if/else，\texttt{demo.py} 还包含循环、异常处理和嵌套结构：}

\vspace{0.05cm}
\begin{columns}[T]
\begin{column}{0.48\textwidth}
{\tiny
\begin{verbatim}
# for 循环
def demo_for_loop(items):
    total = 0
    for item in items:
        total += item
    return total

# while 循环
def demo_while_loop(limit):
    results = []
    counter = 0
    while counter < limit:
        results.append(counter * 2)
        counter += 1
    return results
\end{verbatim}}
\end{column}
\begin{column}{0.48\textwidth}
{\tiny
\begin{verbatim}
# try/except 异常处理
def demo_exception_handling(dividend, divisor):
    try:
        result = dividend / divisor
    except ZeroDivisionError:
        result = 0.0
    else:
        result = result * 1.0
    finally:
        pass
    return result

# 嵌套结构（条件 + 循环）
def demo_nested(data, threshold):
    results = []
    if threshold > 0:
        for item in data:
            count = 0
            while count < threshold:
                results.append(item-count)
                count += 1
    else:
        results = data.copy()
    return results
\end{verbatim}}
\end{column}
\end{columns}
\end{frame}

%------------------------------------------------

\begin{frame}[fragile]
\frametitle{端到端流程（1/3）：静态分析 \texttt{cfg\_builder}}
{\scriptsize 运行 \texttt{python -m pathtracer.cli pathtracer/demo.py} 时，第一步——静态分析：}

\vspace{0.05cm}
{\tiny
\begin{verbatim}
# Step 1: 调用 CFGBuilder 解析 demo.py，找出所有分支点
cfg = CFGBuilder()
all_branches = cfg.build_from_file("pathtracer/demo.py")

# 内部过程：
#   ast.parse(源码)  →  构建 AST 树
#   ast.walk(tree)    →  遍历所有节点
#   对每个节点检查 isinstance(node, ast.If / ast.For / ast.While / ...)

# 结果：找到 15 个分支点
#   Line 32: IF     ← demo_conditionals 里的 if value < 0
#   Line 34: IF     ← elif value == 0  (elif 也是 ast.If)
#   Line 36: IF     ← elif value <= 10
#   Line 38: IF     ← elif value <= 50
#   Line 44: FOR    ← demo_for_loop 里的 for item in items
#   Line 59: WHILE  ← demo_while_loop 里的 while counter < limit
#   Line 67: TRY    ← demo_exception_handling 里的 try
#   Line 69: EXCEPT ← except ZeroDivisionError
#   Line 98: IF     ← demo_nested 里的 if threshold > 0
#   Line 99: FOR    ← for item in data
#   ... 共 15 个
\end{verbatim}}
\end{frame}

%------------------------------------------------

\begin{frame}[fragile]
\frametitle{端到端流程（2/3）：动态追踪 \texttt{tracer}}
{\scriptsize 第二步——用 \texttt{sys.settrace()} 追踪 demo.py 的实际执行过程：}

\vspace{0.05cm}
{\tiny
\begin{verbatim}
# Step 2: 创建追踪器，执行 demo.py 中的所有函数
tracer = PathTracer().trace_file("pathtracer/demo.py")
with tracer:                           # 进入时调用 sys.settrace(callback)
    # 以下是 demo.py 中 main() 实际调用的代码：
    demo_conditionals(25)               # → 走 elif value<=50，"medium"
    demo_conditionals(-5)               # → 走 if value<0，"negative"
    demo_conditionals(0)                # → 走 elif value==0，"zero"
    demo_for_loop([1, 2, 3, 4, 5])     # → for 循环体执行 5 次
    demo_while_loop(5)                  # → while 循环体执行 5 次
    demo_exception_handling(10, 2)      # → try 成功，走 try 分支
    demo_exception_handling(10, 0)      # → 触发 ZeroDivisionError，走 except
    demo_nested([3, 7, 2], 5)          # → if+for+while 嵌套全走
# 退出时调用 sys.settrace(None)，停止追踪

# tracer 在后台记录了每一行的执行：
#   executed_lines = {32,33,34,35,36,37,38,39,44,45,46,47, ...}
#   function_calls = [demo_conditionals(25)→"medium",
#                     demo_conditionals(-5)→"negative", ...]
\end{verbatim}}
\end{frame}

%------------------------------------------------

\begin{frame}[fragile]
\frametitle{端到端流程（3/3）：覆盖率计算 \texttt{reporter}}
{\scriptsize 第三步——对比静态分析的 15 个分支与动态追踪的已执行行号，生成报告：}

\vspace{0.05cm}
{\tiny
\begin{verbatim}
# Step 3: 计算覆盖率
reporter = CoverageReporter()
report = reporter.analyze("pathtracer/demo.py", tracer)

# analyze() 内部：
#   1. cfg_builder.build_from_file() → 15 个分支点
#   2. tracer.get_executed_lines()   → 已执行行号集合 {32,33,34,...}
#   3. 逐分支检查：branch.line_no in executed_lines ?
#      Line 32 (IF):     32 in set? → True  ✓
#      Line 34 (IF):     34 in set? → True  ✓
#      Line 69 (EXCEPT): 69 in set? → True  ✓  (被 (10,0) 触发)
#      ...
#   4. 统计：12 个 True / 15 总共 = 80.0%

# 输出报告：
print(reporter.format_report(report))      # 文本格式 → 终端打印
reporter.format_report_html(report)       # HTML格式 → 源码逐行高亮
reporter.format_report_json(report)       # JSON格式 → 供程序处理
\end{verbatim}}
\end{frame}

%------------------------------------------------

\begin{frame}[fragile]
\frametitle{运行结果：文本报告}
{\footnotesize 运行 \texttt{demo.py} 时，里面\textbf{所有函数都会被执行}（conditionals、for\_loop、while\_loop、exception\_handling、nested……）。报告统计的是\textbf{整个文件}的覆盖情况：}

\vspace{0.05cm}
{\tiny
\begin{verbatim}
============================================================
EXECUTION PATH COVERAGE REPORT
============================================================

Total Branches:   15          ← 来自所有函数的分支点总和
Executed Branches: 12
Coverage: 80.0%

Branches by Type:             ← 按类型分组统计
  if         : 6              ← conditionals 里的 if/elif + nested 里的 if
  for        : 4              ← for_loop + nested 里的 for
  while      : 3              ← while_loop + nested 里的 while
  try        : 1              ← exception_handling 里的 try
  except     : 1              ← exception_handling 里的 except

============================================================
EXECUTED BRANCHES:
------------------------------------------------------------
  demo.py:
    Line   32: if         [executed]
    Line   34: if         [executed]
    Line   44: for        [executed]
    ...
\end{verbatim}}

\vspace{0.05cm}
{\footnotesize 80\% 覆盖率意味着有 3 个分支未被执行。下面单独分析 \texttt{demo\_conditionals} 这个函数，看看为什么不同输入会导致不同覆盖率。}
\end{frame}

%------------------------------------------------

\begin{frame}
\frametitle{聚焦分析：\texttt{demo\_conditionals} 的覆盖对比}
{\footnotesize 单独看 \texttt{demo\_conditionals(value)} 函数：它有 5 个分支（value<0 / ==0 / <=10 / <=50 / else）。
不同的 \texttt{value} 输入让程序走不同的路径，覆盖不同的分支：}

\vspace{0.1cm}
\begin{table}
\centering
{\footnotesize
\begin{tabular}{lcccl}
\toprule
\textbf{测试输入} & \textbf{走了哪条路} & \textbf{覆盖} & \textbf{未覆盖} & \textbf{覆盖率} \\
\midrule
\texttt{demo\_conditionals(-5)} & value<0 $\rightarrow$ "negative" & 1/5 & 4条路径 & 20\% \\
\texttt{demo\_conditionals(0)} & value==0 $\rightarrow$ "zero" & 2/5 & 3条路径 & 40\% \\
\texttt{demo\_conditionals(25)} & value<=50 $\rightarrow$ "medium" & 4/5 & else分支 & 80\% \\
\textbf{综合以上三次调用} & \textbf{覆盖全部5条路} & \textbf{5/5} & \textbf{无} & \textbf{100\%} \\
\bottomrule
\end{tabular}}
\end{table}

\vspace{0.1cm}
{\footnotesize
\begin{block}{结论}
单一测试输入只能覆盖部分路径。要达到高覆盖率，需要\textbf{设计多组不同的测试输入}，覆盖所有分支。\\[2pt]
这正体现了路径分析在软件测试中的价值：\textbf{量化测试充分性}，指导测试用例的补充。
\end{block}}
\end{frame}

%========================================================================================
%	第五部分：总结
%========================================================================================

\section{总结}

\begin{frame}
\frametitle{项目总结}
\begin{block}{做了什么}
编写了 \textbf{PathTracer} 工具，能够自动追踪 Python 程序的执行路径并生成覆盖率报告。
\end{block}

\vspace{0.3cm}
\begin{block}{核心技术}
\begin{enumerate}
    \item \textbf{静态分析}：用 \texttt{ast} 模块解析源代码，识别所有分支点
    \item \textbf{动态追踪}：用 \texttt{sys.settrace()} 在运行时记录执行路径
    \item \textbf{覆盖率计算}：对比静态分支与动态执行结果，计算覆盖率百分比
\end{enumerate}
\end{block}

\vspace{0.3cm}
\begin{block}{输出成果}
\begin{itemize}
    \item 文本报告（终端直接查看）
    \item HTML 可视化报告（源码逐行高亮 + 函数调用树）
    \item JSON 格式（供自动化工具使用）
\end{itemize}
\end{block}
\end{frame}

%------------------------------------------------

\begin{frame}
\frametitle{覆盖的程序结构}
\begin{table}
\centering
\begin{tabular}{lll}
\toprule
\textbf{结构类型} & \textbf{关键词} & \textbf{支持情况} \\
\midrule
顺序结构 & 逐行执行 & \checkmark~行级追踪 \\
选择结构 & \texttt{if / elif / else} & \checkmark~分支覆盖 \\
循环结构 & \texttt{for / while} & \checkmark~进入/跳过 \\
异常处理 & \texttt{try / except} & \checkmark~正常/异常路径 \\
函数调用 & \texttt{def / return} & \checkmark~调用链+参数 \\
嵌套结构 & 组合使用 & \checkmark~递归追踪 \\
\bottomrule
\end{tabular}
\end{table}
\end{frame}

%------------------------------------------------

\begin{frame}
\Huge{\centerline{谢谢！}}
\end{frame}

%----------------------------------------------------------------------------------------

\end{document}
